\chapter{Техніки ухилення, обфускації та антианалізу}
\chaptershort{Техніки ухилення та обфускації}
\label{chap:ch08}

Ефективність методів виявлення й аналізу (розділи~\ref{chap:ch06}--\ref{chap:ch07}) викликала розвиток технік \term{ухилення} --- засобів, якими автори ЗПК приховують справжній код і поведінку від статичного й динамічного аналізу та від сигнатурних і поведінкових детекторів. Ці техніки --- не самостійна функція ЗПК, а оболонка, що захищає корисне навантаження; проте їх розуміння критичне для захисника, бо саме вони визначають межі аналізу (розділ~\ref{chap:ch03}) і диктують вибір методів протидії. Парадоксально, техніки захисту коду (обфускація, пакування) походять із легітимної сфери --- захисту інтелектуальної власності й протидії зворотному аналізу комерційного ПЗ, --- але застосовуються й зловмисниками. Розділ розглядає таксономію ухилення, пакування й шифрування, обфускацію коду, поліморфізм і метаморфізм, антианалізні прийоми, а також методи деобфускації й розпакування, і завершується кейсом підготовки обфускованого зразка до аналізу. Виклад зосереджено на \emph{розпізнаванні} й \emph{подоланні} цих технік, а не на їх створенні.

\section{Таксономія технік ухилення}\label{sec:ch08-taxonomy}

Техніки ухилення класифікують за тим, \emph{що} саме вони мають обманути. \emph{Ухилення від статичного аналізу}: пакування, шифрування, обфускація --- приховують код на диску. \emph{Ухилення від сигнатур}: поліморфізм, метаморфізм --- змінюють байтове подання, зберігаючи функцію. \emph{Ухилення від динамічного аналізу}: антисендбокс, антиВМ, антизневадження, логічні бомби (розділ~\ref{chap:ch07}) --- приховують поведінку за умов спостереження. \emph{Ухилення від виявлення в системі}: руткіт-приховування (розділ~\ref{chap:ch05}), living-off-the-land (розділ~\ref{chap:ch13}) --- маскують присутність під легітимну активність. Узагальнену таксономію наведено в табл.~\thechapter.1.

\tabhead{Таблиця~\thechapter.1 --- Таксономія технік ухилення за ціллю}

\begin{center}\footnotesize
\begin{tabularx}{\linewidth}{L{26mm} Y L{20mm}}
\toprule
Ціль обману & Техніки & Контрзахід \\
\midrule
Статичний аналіз & пакування, шифрування, обфускація & розпакування, деобфускація \\
Сигнатури & поліморфізм, метаморфізм & поведінка, семантика \\
Динамічний аналіз & антиВМ, антизневадження, бомби & прозорі пісочниці \\
Виявлення в системі & руткіти, LotL & аналіз пам'яті, аномалії \\
\bottomrule
\end{tabularx}
\end{center}

Ключове спостереження: жодна техніка ухилення не є абсолютною. Пакування долається розпакуванням, обфускація --- деобфускацією, антиВМ --- прозорими пісочницями. Проте кожна підвищує \emph{вартість} аналізу й знижує ефективність дешевих автоматизованих методів, змушуючи захисника витрачати дорогі ресурси. Тому мета ухилення --- не унеможливити аналіз (це недосяжно для виконуваного коду, що мусить зрештою розгорнутися, аби працювати), а зробити його настільки дорогим, щоб відстрочити виявлення й подовжити час присутності (dwell time, розділ~\ref{chap:ch13}).

\section{Пакування та шифрування}\label{sec:ch08-packing}

\term{Пакувальник} (packer) --- програма, що перетворює виконуваний файл на стиснений і/або зашифрований образ, доповнений \emph{розпакувальником} (stub), який під час запуску відновлює оригінальний код у пам'яті й передає йому керування. Первісне призначення --- стиснення й захист від зворотного аналізу; для ЗПК --- приховування коду від статичного аналізу й сигнатур. Масштабне дослідження пакувальників~\cite{ch08:ugarte} показало широкий спектр складності --- від простого стиснення до багатошарового шифрування з віртуалізацією. Споріднена техніка --- \term{криптор} (crypter), що робить наголос на шифруванні для ухилення від детекторів.

\begin{tikzfig}
\begin{tikzpicture}[font=\scriptsize,
  b/.style={draw,rounded corners=2pt,minimum height=6mm,text width=24mm,align=center,font=\tiny},
  >={Stealth[length=1.5mm]}]
\node[b] (pk) at (0,0) {запакований файл (stub + образ)};
\node[b] (run) at (0,-1.3) {запуск stub};
\node[b] (unp) at (0,-2.6) {розпакування в пам'ять};
\node[b] (oep) at (0,-3.9) {перехід до OEP (справжній код)};
\draw[->] (pk) -- (run); \draw[->] (run) -- (unp); \draw[->] (unp) -- (oep);
\node[b,fill=black!6] (dump) at (4,-2.6) {дамп пам'яті $\to$ статичний аналіз};
\draw[->,dashed] (oep.east) to[out=0,in=-90] (dump.south);
\end{tikzpicture}
\figcap{Рис.~\thechapter.1 --- Життєвий цикл запакованого зразка й точка перехоплення для розпакування (дамп після досягнення OEP)}
\end{tikzfig}

Пакування майже повністю знищує цінність байтових сигнатур: один і той самий код під різними пакувальниками (або з різними ключами) дає зовсім різні файли. Ознакою пакування слугує висока ентропія секцій~\eqref{eq:ch03-entropy}, невідповідність між оголошеним та фактичним розміром, аномальна таблиця імпортів (часто лише кілька функцій для динамічного завантаження решти), нетипові імена секцій. Проте пакування безсиле проти динамічного аналізу й аналізу пам'яті (розділ~\ref{chap:ch07}): щоб виконатися, код \emph{мусить} розпакуватися, після чого його можна зняти дампом. Ця фундаментальна властивість --- неможливість виконати прихований код, не розкривши його, --- є ключем до подолання пакування.

Складність пакувальників утворює спектр~\cite{ch08:ugarte}. Найпростіші виконують одношарове стиснення зі статичним розпакувальником, який легко розпізнати й автоматично обернути. Складніші додають кілька шарів шифрування, перевірки цілісності (щоб аналітик не міг пропатчити stub), антианалізні перевірки в самому розпакувальнику та відкладене розгортання коду частинами (щоб у пам'яті ніколи не був присутній увесь код одразу). Найскладніші поєднують пакування з віртуалізацією самого розпакувальника. Проте навіть вони підкоряються базовому обмеженню: будь-яка інструкція, що виконується процесором, у момент виконання присутня у відкритому вигляді в пам'яті. Тому стратегія захисника --- не намагатися статично «прочитати» запакований файл, а дозволити йому розгорнутися під контролем і зафіксувати відкритий код. Різниця між пакувальниками зводиться до того, \emph{коли} й \emph{наскільки повно} код розкривається, а не до того, чи розкривається взагалі.

Слід відрізняти \emph{стандартні} пакувальники загального призначення від \emph{кастомних}. Стандартні (стискачі та протектори, доступні публічно) мають відомі сигнатури stub-коду й типові розпакувальники, тому часто розпізнаються й автоматично обертаються інструментами. Саме тому серйозні зловмисники застосовують власні або суттєво модифіковані пакувальники, унікальні для кампанії: вони позбавлені відомих сигнатур і потребують ручного аналізу. Парадоксально, сам факт застосування рідкісного кастомного пакувальника стає ознакою: легітимне ПЗ зазвичай користується відомими комерційними протекторами, тож нестандартна оболонка підвищує підозрілість зразка. Це приклад загального принципу --- зусилля з приховування самі лишають слід, який захисник може використати як ознаку.

\section{Обфускація коду}\label{sec:ch08-obfuscation}

\term{Обфускація} --- перетворення коду, що зберігає його функцію, але ускладнює розуміння~\cite{ch08:collberg}. На відміну від пакування (приховує код цілком до розпакування), обфускація лишає код виконуваним і видимим, але заплутаним. Основні прийоми: \emph{непрозорі предикати} (умови, значення яких завжди відоме автору, але важко обчислюване аналізатором, --- створюють фіктивні гілки); \emph{сплощення потоку керування} (control-flow flattening --- заміна природної структури CFG на диспетчер із великим перемикачем, що приховує справжній порядок блоків); \emph{підстановка еквівалентних інструкцій} та вставлення мертвого коду; \emph{віртуалізація коду} (переклад коду на власну байт-кодову машину з інтерпретатором --- найстійкіший і найдорожчий прийом). Інструменти на кшталт Obfuscator-LLVM~\cite{ch08:junod} автоматизують ці перетворення.

\begin{tikzfig}
\begin{tikzpicture}[font=\scriptsize,
  n/.style={draw,rounded corners=2pt,minimum height=4.5mm,text width=11mm,align=center,font=\tiny,inner sep=1.5pt},
  d/.style={draw,diamond,aspect=2,inner sep=1pt,font=\tiny},
  >={Stealth[length=1.3mm]}]
% left: natural CFG
\node[n] (a1) at (0,0) {A}; \node[n] (a2) at (0,-0.9) {B};
\node[n] (a3) at (0,-1.8) {C}; \node[n] (a4) at (0,-2.7) {D};
\draw[->] (a1)--(a2); \draw[->] (a2)--(a3); \draw[->] (a3)--(a4);
\node[font=\tiny] at (0,-3.4) {природний};
% right: flattened
\node[d] (dsp) at (3.6,0) {disp};
\node[n] (b1) at (2.4,-1.3) {A}; \node[n] (b2) at (3.2,-1.3) {B};
\node[n] (b3) at (4.0,-1.3) {C}; \node[n] (b4) at (4.8,-1.3) {D};
\draw[->] (dsp)--(b1); \draw[->] (dsp)--(b2); \draw[->] (dsp)--(b3); \draw[->] (dsp)--(b4);
\draw[->] (b1) to[out=-90,in=-90] (dsp.south); \draw[->] (b4) to[out=-90,in=-90] (dsp.south);
\node[font=\tiny] at (3.6,-2.6) {сплощений (диспетчер)};
\end{tikzpicture}
\figcap{Рис.~\thechapter.2 --- Сплощення потоку керування: природний CFG (ліворуч) замінюється диспетчером із перемикачем (праворуч), що приховує справжній порядок блоків}
\end{tikzfig}

Обфускація безпосередньо б'є по аналізу коду (розділ~\ref{chap:ch06}): сплощення руйнує зрозумілий CFG, непрозорі предикати породжують фіктивні шляхи, віртуалізація ховає логіку за інтерпретатором. Розгляньмо прийоми докладніше. \emph{Непрозорий предикат} --- це вираз, значення якого автор знає наперед (наприклад, тотожність на кшталт $7x^2 - 1 \ne y^2$, що істинна для всіх цілих), але доведення якого вимагає від аналізатора нетривіальних міркувань; вставлений як умова розгалуження, він створює гілку, що насправді ніколи (або завжди) не виконується, роздуваючи CFG фіктивними шляхами. \emph{Сплощення потоку керування} (рис.~\thechapter.2) розбиває функцію на базові блоки, нумерує їх і підпорядковує єдиному диспетчеру-перемикачу, який за значенням керованої змінної обирає наступний блок; природна послідовність зникає, і відновити її можна лише, простеживши значення керованої змінної. \emph{Віртуалізація} замінює машинні інструкції байт-кодом власної віртуальної машини, вбудованої у зразок; аналітик мусить спершу зрозуміти саму цю машину, перш ніж декодувати логіку, --- що й робить віртуалізацію найдорожчим для подолання прийомом. Проте й вона має теоретичні межі: доведено, що «ідеальна» обфускація неможлива в загальному випадку, а практична обфускація лише підвищує вартість аналізу~\cite{ch08:oreans}. Семантичний і динамічний аналіз, що працюють з ефектами коду, а не з його записом, частково нейтралізують обфускацію, бо яким би заплутаним не був код, його спостережувана поведінка лишається тією самою.

\section{Поліморфізм та метаморфізм}\label{sec:ch08-poly}

\term{Поліморфізм} --- здатність коду відтворювати себе в різних байтових формах зі збереженням функції, класично через шифрування тіла змінним ключем і мутацію розшифровувача~\cite{ch08:szor}. \term{Метаморфізм} --- глибша техніка: мутація \emph{самого} коду (перестановка інструкцій, заміна еквівалентними, зміна регістрів) без шифрування, так що кожне покоління має інший код тієї самої функції. Обидві техніки націлені на сигнатури: якщо кожен екземпляр унікальний на рівні байтів, байтова сигнатура безсила.

Ефективність ухилення від сигнатур зручно описати комбінаторно. Якщо генератор варіантів має $m$ незалежних точок мутації, кожна з $r$ рівноможливими варіантами, кількість різних байтових форм

\begin{equation}
N = r^{m},
\label{eq:ch08-variants}
\end{equation}

зростає експоненційно, роблячи перелік усіх форм однією байтовою сигнатурою неможливим. Формула~\eqref{eq:ch08-variants} пояснює, чому проти поліморфізму й метаморфізму ефективні лише методи, інваріантні до таких мутацій: виявлення за розшифрованим тілом у пам'яті, поведінкові сигнатури (розділ~\ref{chap:ch07}), семантичний аналіз (розділ~\ref{chap:ch06}), статистичні ознаки (розділ~\ref{chap:ch09}). Саме перехід від байтового до поведінкового й семантичного рівня є системною відповіддю на поліморфізм.

Історично поліморфізм і метаморфізм --- вершина еволюції самозахисту вірусної епохи~\cite{ch08:szor,ch08:you}, коли виявлення спиралося майже виключно на байтові сигнатури. Поліморфний генератор шифрує незмінне тіло вірусу змінним ключем і щоразу по-новому будує розшифровувач (переставляючи інструкції, вставляючи «сміття», змінюючи регістри), тож єдиним спільним інваріантом лишається поведінка розшифрованого тіла в пам'яті. Відповіддю антивірусної галузі стала емуляція: зразок виконують у вбудованому емуляторі, доки він сам не розшифрує тіло, після чого застосовують сигнатуру до відкритого коду --- рання форма динамічного подолання ухилення. Метаморфізм пішов далі, відмовившись від сталого тіла взагалі: мутує весь код, і спільного байтового інваріанта не лишається зовсім. Це остаточно засвідчило вичерпність суто байтового підходу й прискорило розвиток поведінкових, емуляційних і статистичних методів, які й становлять сучасну основу виявлення.

\section{Антианалізні прийоми}\label{sec:ch08-antianalysis}

Антианаліз протидіє безпосередньо роботі аналітика та інструментів~\cite{ch08:branco}. \emph{Антизневадження} (anti-debugging): виявлення зневаджувача через прапорці системи, вимірювання часу виконання (у зневаджувачі код повільніший), навмисні винятки, перевірку контрольних точок. \emph{Антидизасемблювання}: вставлення байтів, що збивають лінійний прохід дизасемблера (розділ~\ref{chap:ch06}), перекриття інструкцій, непрямі переходи. \emph{АнтиВМ та антисендбокс} (розділ~\ref{chap:ch07}): виявлення віртуалізації та інструментування. \emph{Антиемуляція}: використання рідкісних або недокументованих інструкцій, що неточно відтворюються емулятором.

Окрему роль відіграє \emph{часовий} антианаліз. Виконання під зневаджувачем або в емуляторі значно повільніше за виконання на реальному процесорі, тому зразок вимірює тривалість певної послідовності операцій і, виявивши аномальну затримку, робить висновок про спостереження. Споріднений прийом --- вимірювання часу між двома точками з очікуванням, що між ними не було зовнішнього втручання: контрольна точка зневаджувача або перехоплення розтягує цей інтервал. Контрзаходи --- нормалізація часу (підміна показань таймерів так, щоб вони відповідали «реальній» швидкості) та перехід на апаратне трасування, невидиме для зразка. Часовий антианаліз показово ілюструє загальний принцип: будь-яка спроба ухилення потребує \emph{спостережуваної} відмінності середовищ, і усунення цієї відмінності знімає відповідну перевірку.

Спільна риса антианалізу --- опора на \emph{розрізнюваність} середовища аналізу й реальної системи. Доки аналітичне середовище має помітні відмінності (прапорці, часові аномалії, артефакти ВМ), зразок може їх виявити. Контрзахід --- підвищення прозорості: приховування зневаджувача, нормалізація часу, bare-metal-аналіз (розділ~\ref{chap:ch07}), а також виявлення самих антианалізних перевірок як підозрілої ознаки --- наявність численних перевірок середовища сама свідчить про зловмисність. Крім того, антианалізний код підвищує вартість, але не робить аналіз неможливим: досвідчений аналітик обходить перевірки, патчить гілки або працює на рівні, недосяжному для перевірок зразка.

\section{Деобфускація та розпакування}\label{sec:ch08-deobf}

Подолання ухилення --- окрема дисципліна. \term{Розпакування} спирається на фундаментальну властивість: код мусить розгорнутися, аби виконатися. \emph{Загальне (універсальне) розпакування}~\cite{ch08:royal,ch08:guo} не залежить від конкретного пакувальника: зразок виконують, відстежуючи пам'ять, і фіксують момент, коли керування передається в раніше записану (а отже, розпаковану) ділянку. Інваріант «записано-потім-виконано» (write-then-execute) є ознакою завершення розпакування: якщо байт за адресою спершу записано під час виконання, а потім з нього вибирають інструкцію, це майже напевно розпакований код. Формально ознаку досягнення точки входу оригіналу можна виразити умовою

\begin{equation}
\exists a:\; W(a) \prec X(a),
\label{eq:ch08-wx}
\end{equation}

де $W(a)$ --- подія запису за адресою $a$, $X(a)$ --- подія виконання коду за тією самою адресою, а $\prec$ --- часове передування. Виявивши таку подію, аналізатор знімає дамп пам'яті й продовжує статичний аналіз уже відкритого коду.

\term{Деобфускація}~\cite{ch08:yadegari} відновлює зрозумілий код з обфускованого: спрощення непрозорих предикатів (через виявлення завжди-істинних умов), відновлення природного CFG зі сплощеного, усунення мертвого коду, девіртуалізація. Загальні підходи спираються на динамічний аналіз потоку даних і символьне виконання, що працюють із семантикою, а не з заплутаним записом. Вітчизняні дослідження систематизують методи обфускації та підходи до їх подолання~\cite{ch08:obfusc-ua}, а також виявлення запакованих зразків за ентропією та структурними ознаками~\cite{ch08:packer-ua}.

\section{Кейс: підготовка обфускованого зразка до аналізу}\label{sec:ch08-case}

Розглянемо послідовність подолання ухилення для аналізу запакованого й обфускованого зразка (захисна задача).

\emph{Крок 1. Розпізнавання захисту.} За ентропією~\eqref{eq:ch03-entropy}, структурою секцій та імпортів визначаємо наявність пакування й, за можливості, його тип. Висока ентропія й бідна таблиця імпортів --- типові ознаки.

\emph{Крок 2. Розпакування.} Застосовуємо загальне розпакування: виконуємо зразок у контрольованому середовищі, відстежуючи пам'ять за інваріантом~\eqref{eq:ch08-wx}. Досягнувши точки входу оригіналу, знімаємо дамп і відновлюємо таблицю імпортів.

\emph{Крок 3. Оцінювання обфускації.} На розпакованому коді оцінюємо рівень обфускації: сплощення CFG, непрозорі предикати, ознаки віртуалізації. Це визначає потрібну глибину деобфускації.

\emph{Крок 4. Деобфускація.} Спрощуємо непрозорі предикати, відновлюємо природний CFG, усуваємо мертвий код, за потреби девіртуалізуємо. Для стійких до статики прийомів застосовуємо динамічний і семантичний аналіз, що працюють з поведінкою.

\emph{Крок 5. Аналіз і сигнатури.} На очищеному коді проводимо повноцінний статичний аналіз (розділ~\ref{chap:ch06}), виділяємо стійкі семантичні ознаки та будуємо правила виявлення, інваріантні до повторного пакування (орієнтовані на поведінку й семантику, а не на байти оболонки).

Кейс ілюструє стратегічний принцип: ухилення підвищує вартість, але не скасовує аналізу, бо виконуваний код зрештою мусить розгорнути свою справжню логіку. Захисник переносить виявлення на рівні, інваріантні до ухилення, --- поведінку, семантику, аномалії, --- замість того щоб змагатися з оболонкою на її ж полі байтових форм.

\section*{Висновки до розділу}
\addcontentsline{toc}{section}{Висновки до розділу}

Техніки ухилення --- пакування, шифрування, обфускація, поліморфізм, метаморфізм, антианаліз --- є захисною оболонкою ЗПК, спрямованою проти конкретних методів виявлення й аналізу. Їх класифікують за ціллю обману: статичний аналіз, сигнатури, динамічний аналіз, виявлення в системі (табл.~\thechapter.1). Пакування приховує код на диску, знищуючи байтові сигнатури, але безсиле проти динамічного аналізу й аналізу пам'яті, бо код мусить розпакуватися для виконання. Обфускація лишає код видимим, але заплутаним, б'ючи по аналізу коду, проте має теоретичні межі й нейтралізується семантичним і динамічним аналізом. Поліморфізм і метаморфізм породжують експоненційну кількість байтових форм~\eqref{eq:ch08-variants}, роблячи байтові сигнатури безсилими й змушуючи переходити до поведінкового й семантичного рівня. Антианаліз спирається на розрізнюваність середовища аналізу й долається підвищенням прозорості. Деобфускація й загальне розпакування спираються на фундаментальну властивість виконуваного коду --- необхідність розгорнутися~\eqref{eq:ch08-wx}. Головний висновок: ухилення підвищує вартість аналізу, але не скасовує його; системна відповідь захисника --- перенесення виявлення на рівні, інваріантні до ухилення (поведінка, семантика, аномалії), замість змагання з оболонкою на рівні байтів.

\selfcheck
\begin{enumerate}
\item За якою ознакою класифікують техніки ухилення (табл.~\thechapter.1)? Чому жодна з них не є абсолютною?
\item Яка справжня мета технік ухилення, якщо аналіз виконуваного коду принципово неможливо унеможливити?
\item Що таке пакувальник і криптор? Чому пакування знищує цінність байтових сигнатур?
\item Назвіть ознаки пакування. Чому пакування безсиле проти динамічного аналізу й аналізу пам'яті?
\item Що таке обфускація? Опишіть непрозорі предикати, сплощення CFG та віртуалізацію коду.
\item Чому обфускація має теоретичні межі? Як семантичний і динамічний аналіз її нейтралізують?
\item Порівняйте поліморфізм і метаморфізм. Проти якого методу виявлення вони спрямовані?
\item За формулою~\eqref{eq:ch08-variants} поясніть, чому байтові сигнатури безсилі проти поліморфізму. Які методи ефективні натомість?
\item Опишіть антизневадження, антидизасемблювання та антиемуляцію. На чому ґрунтується їх спільна риса?
\item Чому наявність численних перевірок середовища сама є ознакою зловмисності?
\item Поясніть принцип загального розпакування та інваріант «записано-потім-виконано»~\eqref{eq:ch08-wx}.
\item Що таке деобфускація? Які підходи застосовують для відновлення зрозумілого коду?
\item Опишіть послідовність підготовки обфускованого зразка до аналізу (кейс підрозділу~\ref{sec:ch08-case}).
\item Дослідницьке: оцініть, наскільки різні пакувальники підвищують вартість аналізу й знижують ефективність сигнатурного виявлення на наборі зразків.
\item Дослідницьке: порівняйте стійкість поведінкових та семантичних ознак до повторного пакування й обфускації того самого корисного навантаження.
\end{enumerate}
